\section{Background and Requirements}
\label{sec:background}

This chapter introduces the electrical concepts a reader needs to
follow the firmware design --- solar-panel power transfer, MPPT,
buck conversion, complementary PWM, eFuses, framed serial
communication, and the SAMD21 platform --- and the requirements
inherited from the CHESS mission document.

\subsection{The EPS in one diagram}

The shortest useful model of the EPS is a chain:

\begin{center}
\texttt{solar arrays + battery $\rightarrow$ PCU conditions and
protects energy $\rightarrow$ PDU distributes protected rails
$\rightarrow$ MCU measures, decides, commands, reports.}
\end{center}

The PCU is responsible for the energy-conversion physics. The PDU
distributes the conditioned rails to the loads (OBC, payload, radio,
attitude control, ADCS) through eFuses. The MCU --- the subject of
this report --- is the supervisor that ties the two halves
together and speaks to the OBC.

\subsection{Operating modes of the PCU}
\label{sec:pcu-modes}

The PCU has four normal operating modes, each driven by the current
combination of solar availability and battery state of charge.

\begin{description}
\item[\texttt{MPPT\_CHARGE}] The Sun is available and the battery is
not yet full. The PCU draws as much energy as possible from the
panels and routes it to the battery. This is where the MPPT
algorithm is active.

\item[\texttt{CV\_FLOAT}] The battery is at or near its maximum
safe voltage. The PCU stops trying to maximise panel power and
instead regulates the charging rail to a constant voltage, letting
charge taper off naturally.

\item[\texttt{SA\_LOAD\_FOLLOW}] The battery is fully charged and
the Sun is still available. The PCU runs the array only hard enough
to cover the spacecraft's instantaneous load, avoiding overcharge.

\item[\texttt{BATTERY\_DISCHARGE}] No usable sunlight (the satellite
is in eclipse, or the panels are still stowed). The panel rail is
opened and the battery powers everything.
\end{description}

Above these four sits the \texttt{SAFE\_MODE} state, an emergency
override entered when the battery voltage drops below a critical
threshold, the battery temperature is outside its safe operating
range, or the OBC has been silent for longer than the heartbeat
timeout. Safe mode has four sub-states selected by the OBC
(\texttt{DETUMBLING}, \texttt{CHARGING}, \texttt{COMMUNICATION},
\texttt{REBOOT}); each sub-state defines the load mask that must
stay powered while the satellite is in survival configuration.

The transition rules and the threshold definitions are documented
in \texttt{state-transitions.md} inside the project repository, and
implemented in pure C as a single dispatcher function described in
Chapter~\ref{sec:architecture}.

\subsection{Maximum Power Point Tracking}

A solar panel does not behave as a constant voltage source. Its
current--voltage (I--V) curve has a single point of maximum power
output that depends on irradiance, temperature, and load. The job
of an MPPT algorithm is to actively steer the converter's operating
point onto that maximum.

The mission document selects Incremental Conductance as the MPPT
algorithm of choice. Incremental Conductance compares the
instantaneous conductance $I/V$ to the incremental conductance
$\mathrm{d}I/\mathrm{d}V$: at the maximum power point the sum of
the two is zero. Adjusting the buck converter's duty cycle moves
the operating point along the I--V curve, and the algorithm walks
the duty cycle up or down to drive that sum toward zero. Compared
to Perturb-and-Observe, Incremental Conductance is less prone to
oscillation around the peak and tracks faster under changing
irradiance.

\subsection{The buck converter and complementary PWM}

The PCU's energy-conversion element is a step-down (buck) DC--DC
converter built around an EPC2152 GaN half-bridge. The EPC2152
needs two synchronised gate-drive inputs --- high-side and
low-side --- that are complementary: when one is on, the other must
be off. If both were on simultaneously, even for nanoseconds, the
half-bridge would short the input rail (a condition called
\textit{shoot-through}) and the FETs would be destroyed.

The complementary drive is generated by the SAMD21's TCC0 timer
peripheral. Both gates run at 300\,kHz with the duty cycle
controlled by software. A small \textit{dead-time} interval, during
which both gates are off, is inserted at every transition. The
SAMD21 datasheet specifies the dead-time generator as a hardware
feature of the TCC peripheral: the firmware programs two registers
(\texttt{DTLS} and \texttt{DTHS}) to set the low-side and
high-side dead times in units of the clock period.

For the CHESS hardware the dead time is set to approximately
100\,ns. At 48\,MHz core clock this corresponds to five clock
periods. The value was confirmed visually on an oscilloscope during
the mainboard bring-up.

\subsection{The CHIPS protocol}

CHIPS (CHESS Internal Protocol over Serial) is the framing protocol
used across the satellite for every OBC--subsystem exchange. The
mission document defines its frame structure:

\begin{itemize}
\item Two sync bytes (\texttt{0x7E}) delimit every frame.
\item A one-byte sequence number, a one-byte command/response
byte (with the high bit indicating direction), and a variable-length
payload follow.
\item A two-byte CRC-16/KERMIT checksum is appended before the
trailing sync byte.
\item Byte stuffing (PPP-style, RFC~1662) escapes any payload byte
equal to \texttt{0x7E} or \texttt{0x7D}.
\end{itemize}

The OBC is always the master: it issues commands; each subsystem
acknowledges within one second. If the OBC sends the same sequence
number twice in a row (because its first request timed out), the
subsystem must execute the command only once but reply each time
the request arrives. After 120 seconds without an OBC message, a
subsystem may assume autonomy.

\subsection{The SAMD21 platform}

The MCU is a Microchip SAMD21J17D-MUT on the real PCU board
(64-pin QFN, 128\,kB flash, 16\,kB RAM) and a SAMD21G17D on the
Curiosity Nano dev board used for early bring-up (48-pin TQFP, same
silicon die). Both are Cortex-M0+ cores at up to 48\,MHz. The
choice of bare-metal C without HAL or RTOS is deliberate: the same
flash budget that buys a HAL on a richer chip is needed here for
the application logic itself, and direct register access makes
real-time behaviour straightforward to audit.

The relevant peripherals used by the EPS firmware are TCC0 for the
buck PWM, SERCOM0 for the OBC UART, SERCOM3 for the I\textsuperscript{2}C bus
to the INA226 chips, and the ADC for the analog sensor inputs.

\subsection{Mission requirements traceability}

The detailed requirement set lives in the project repository as
\texttt{docs/requirements.md}. The most relevant requirements for
this report are reproduced in Table~\ref{tab:requirements}; each
maps to a section of the architecture chapter or the test plan.

\begin{table}[ht]
\centering
\caption{Selected mission requirements relevant to this report.}
\label{tab:requirements}
\small
\begin{tabular}{@{}lp{0.66\linewidth}@{}}
\toprule
\textbf{ID} & \textbf{Requirement} \\
\midrule
PCU-1 & The PCU shall use a buck converter driven by complementary
PWM with dead-time. \\
PCU-2 & The buck converter shall switch at 300\,kHz. \\
PCU-3 & The PCU shall run MPPT when solar power is available and
the battery needs charge. \\
PCU-4 to PCU-7 & The PCU shall support the four operating modes
\texttt{MPPT\_CHARGE}, \texttt{CV\_FLOAT},
\texttt{SA\_LOAD\_FOLLOW}, \texttt{BATTERY\_DISCHARGE}. \\
MCU-1 & The MCU shall boot first and manage power before the OBC is
available. \\
MCU-3 to MCU-5 & The MCU shall read PCU/PDU electrical telemetry
via I\textsuperscript{2}C and ADC. \\
MCU-9 & The MCU shall communicate with the OBC over CHIPS. \\
MCU-12 & The MCU shall enter autonomous behaviour after 120\,s of
OBC silence. \\
SAFE-1 to SAFE-8 & The MCU shall report critical conditions, shed
non-essential loads, activate heaters when needed, and not exit
safe mode autonomously. \\
\bottomrule
\end{tabular}
\end{table}

Every requirement is traceable to a section of the implementation
(Chapter~\ref{sec:architecture}) and to a verification step
(Chapter~\ref{sec:hil}).
